% Edit this file directly. It was generated once from paper.tex (commit ad5eb72) by build_p1.py, which is now frozen.
\newcommand{\papertitle}{A log-concavity criterion for finite nonnegative linearization}
\newcommand{\papershort}{A LOG-CONCAVITY CRITERION FOR FINITE LINEARIZATION}
\newcommand{\paperdate}{October 23, 2026}
\documentclass[twoside,10pt]{article}
\usepackage[T1]{fontenc}
\usepackage{lmodern}
\usepackage{amsmath,amssymb,amsthm,amsfonts,mathtools}
\usepackage[letterpaper,total={5.5in,8in},vmarginratio=1:1,hmarginratio=1:1,headsep=21pt]{geometry}
\usepackage{xcolor}
\usepackage{xr-hyper}
\usepackage[breaklinks=true,linktocpage=true]{hyperref}
\hypersetup{colorlinks=true,linkcolor=purple,urlcolor=purple,citecolor=blue!80!black,
pdfauthor={Leslie P. Polzer},pdftitle={\papertitle},
pdfsubject={Finite orthogonal polynomials and nonnegative linearization},
pdfkeywords={orthogonal polynomials, Krawtchouk polynomials, nonnegative linearization, quadratic transformation, hypergroups, Lean}}
\medmuskip=4mu
\thickmuskip=8mu
\setlength{\emergencystretch}{2em}
\numberwithin{equation}{section}
\newtheoremstyle{slbody}{8pt}{8pt}{\normalfont\slshape}{}{\bfseries}{.}{0.5em}
{\thmname{#1}\thmnumber{ #2}\thmnote{ {\normalfont\itshape(\ignorespaces#3\unskip)}}}
\theoremstyle{slbody}
\newtheorem{theorem}{Theorem}[section]
\newtheorem{lemma}[theorem]{Lemma}
\newtheorem{corollary}[theorem]{Corollary}
\newtheorem{proposition}[theorem]{Proposition}
\theoremstyle{definition}
\newtheorem{definition}[theorem]{Definition}
\theoremstyle{remark}
\newtheorem{remark}[theorem]{Remark}
\renewcommand{\qedsymbol}{$\blacksquare$}
\newcommand{\qint}[1]{[#1]_q}
\title{\papertitle}
\author{Leslie P. Polzer}
\newcommand{\authoraffiliation}{Independent Researcher}
\newcommand{\authoremail}{polzer@fastmail.com}
\date{\paperdate}
\makeatletter
\renewcommand{\section}{\@startsection{section}{1}{\z@}{-3.5ex \@plus -1ex \@minus -.2ex}{2.3ex \@plus .2ex}{\normalfont\large\bfseries}}
\renewcommand{\@seccntformat}[1]{\csname the#1\endcsname.\enspace}
\def\tagform@#1{\maketag@@@{(\ignorespaces{\oldstylenums{#1}}\unskip\@@italiccorr)}}
\renewcommand{\eqref}[1]{\textup{{\normalfont(\oldstylenums{\ref{#1}}}\normalfont)}}
\newcommand{\shorttitle}{\papershort}
\newcommand{\shortauthors}{LESLIE P. POLZER}
\newcommand{\ps@bookheader}{%
\renewcommand\@oddfoot{\hfil}\renewcommand\@evenfoot{\hfil}%
\renewcommand\@oddhead{\ifnum\value{page}>1{\small\hfil\shortauthors}\hfil\thepage\else\hfil\fi}%
\renewcommand\@evenhead{\thepage\hfil\small\shorttitle\hfil}}
\pagestyle{bookheader}
\renewcommand{\maketitle}{\begin{center}{\large\bfseries\@title}\\\vskip36pt
{\scshape\@author}\par\medskip{\small\authoraffiliation}\par\smallskip
{\small\texttt{\authoremail}}\par\smallskip{\small\@date}\par\vskip36pt\end{center}}
\makeatother
\renewenvironment{abstract}{\quotation\noindent\small{\bfseries\abstractname.\enspace}}{\endquotation}

\begin{document}
\maketitle

\begin{abstract}
Szwarc asked for criteria, stated in the recurrence coefficients, that give nonnegative linearization (property (P)) for finite systems of orthogonal polynomials such as the Krawtchouk polynomials, and for a criterion that handles his symmetric quadratic system $Q_n$. For a finite stochastic three-term recurrence whose diagonal and downward coefficients have a constant ratio, we prove that property (P) holds for every product, including products beyond the terminal degree, whenever the Newton coefficients of the profile $d_i=b_i+c_i$ are log-concave without internal zeros. The criterion contains the full positive Krawtchouk range and the geometric (dual $q$-Krawtchouk) arrays. Through a quadratic decomposition of symmetric finite systems it gives a parity-split criterion that covers $Q_n$ for every size and every $0<p_0\leq1/2$, which is sharp. For every fixed $0<r<1$, no criterion valid for all sizes contains coefficient balls of a fixed radius around the Krawtchouk arrays. Positivity of the Newton coefficients does not suffice, while log-concavity is not necessary; we pose the exact boundary as a problem. Its theorems, corollaries and propositions are formalized in Lean.
\vskip5pt
\noindent\textbf{Keywords.}\enspace Finite orthogonal polynomials; nonnegative linearization; Krawtchouk polynomials; log-concave sequences; quadratic transformation; polynomial hypergroups.
\vskip5pt
\noindent\textbf{MSC2020 Classification.}\enspace 33C45; 42C05; 05A20.
\end{abstract}

\vskip50pt
\baselineskip=13.5pt
\section{Introduction}\label{sec:introduction}

Nonnegative linearization (property (P)) asks whether the product of two orthogonal polynomials is a nonnegative linear combination of polynomials of the same family. For a finite system $R_0,\ldots,R_N$ the identity is understood on the orthogonality support, so it includes products whose degree sum exceeds $N$. In the normalization $R_i(1)=1$ the coefficients then sum to one and define a finite polynomial hypergroup.

In his survey, Szwarc~\cite[\S5, Problem~2]{szwarc} asks: ``Find any criteria for the property (P) for finite systems of orthogonal polynomials, like the Krawtchouk polynomials.''\footnote{In the author version, \S5, Problem~2, printed p.~125, the displayed weights $\binom Nn p^{N-n}(1-p)^n$ have mean $N(1-p)$, whereas the displayed recurrence gives $K_1(x)=(x-pN)/((1-p)N)$, so that $\int K_1\,d\mu=(1-2p)/(1-p)$. Corollary~\ref{cor:kraw} fixes the convention directly: it uses the displayed weights and the normalization at the point carrying the weight $p^N$.} He then introduces a symmetric system $Q_n$ built from the Krawtchouk recurrence, notes that it has property (P) for $p_0\leq1/2$, and asks for a general criterion that would handle this case. His criteria for infinite systems~\cite[Theorems~3.2 and~3.3]{szwarc} are monotonicity conditions on the recurrence coefficients; for a finite system some of them fail near the terminal index, whose upward coefficient vanishes (Section~\ref{sec:szwarc}). Askey~\cite[Part~II, Theorem~1 and Corollary~1]{askey} gave finite recurrence conditions for products of degree sum at most $N$, through a maximum principle for a hyperbolic difference equation, and Szwarc~\cite[Theorem~1]{szwarc2003} characterized property (P) for finite and infinite systems through an auxiliary discrete hyperbolic problem. Miclo~\cite[Proposition~3.3]{miclo} characterizes the hypergroup property of a finite birth--death chain at a boundary point spectrally. None of these gives a coefficient criterion that covers all products of the finite Krawtchouk systems.

We work with a stochastic three-term recurrence \eqref{eq:recurrence} with diagonal coefficients $b_i$, downward coefficients $c_i$ and profile $d_i=b_i+c_i$. The main result, Theorem~\ref{thm:lc}, concerns arrays with a constant ratio $b_i/c_i$: property (P) holds for every product whenever the Newton coefficients $\nu_k=\Delta^kd(0)$ of the profile are nonnegative and log-concave without internal zeros. This is a set of inequalities on the coefficients; checking it does not involve any product coefficient. It contains the full positive Krawtchouk range (Corollary~\ref{cor:kraw}; property (P) is classical there, see Eagleson~\cite{eagleson} and~\cite[\S5]{szwarc}) and the geometric arrays of Corollary~\ref{cor:additional}, the dual $q$-Krawtchouk systems of Koelink's construction on the Hecke algebra of type~B~\cite{koelink}. Section~\ref{sec:szwarc} answers the second part of the problem: a quadratic decomposition reduces property (P) of a symmetric finite system to two nonsymmetric ones (Theorem~\ref{thm:decomp}), and Theorem~\ref{thm:lc} then gives a parity-split criterion (Corollary~\ref{cor:parity}) that covers $Q_n$ for every $N$ and $0<p_0\leq1/2$; this range is sharp (Theorem~\ref{thm:szwarc}).

Section~\ref{sec:limits} discusses how far criteria of this kind can go. For every fixed $r<1$, no set of arrays with property (P) contains coefficient balls of a radius independent of $N$ around the Krawtchouk arrays (Corollary~\ref{cor:radius-exp}), so a criterion valid for all $N$ must place them near its boundary; the log-concavity condition places them on its edge. Positivity of the Newton coefficients does not suffice, but log-concavity is not necessary; we record exact examples and pose the exact boundary as a problem.

Every theorem, corollary and proposition, and Lemmas~\ref{lem:finite} and~\ref{lem:corner-value}, is formalized in Lean~4 with Mathlib; the lemmas of Section~\ref{sec:lc} are formalized within the proof of Theorem~\ref{thm:lc} (see the section on code and formalization). Two companion papers continue this work: \cite{P2} treats ratios $c_i/d_i$ that increase with the index, and \cite{P3} treats the algebraic geometric-increment criterion, explicit perturbation neighbourhoods and the constraints at the symmetric Krawtchouk endpoint.

\section{Setting}\label{sec:criterion}

For integers $j\geq0$ put $[j]_q=1+q+\cdots+q^{j-1}$, with $[0]_q=0$; in particular $[j]_1=j$.
Consider real coefficients indexed by $0\leq i\leq N$, where $N\geq1$, satisfying
\begin{equation}\label{eq:basic}
\begin{gathered}
a_0=1,\quad b_0=c_0=0,\quad a_N=0,\quad a_i+b_i+c_i=1,\\
a_i>0\ (i<N),\qquad c_i>0\ (i>0).
\end{gathered}
\end{equation}
Set $R_{-1}=0$, $R_0=1$, and determine $R_1,\ldots,R_N$ by
\begin{equation}\label{eq:recurrence}
xR_i=a_iR_{i+1}+b_iR_i+c_iR_{i-1}\qquad(0\leq i<N).
\end{equation}
At $i=0$, \eqref{eq:recurrence} gives $R_1=x$; induction at $x=1$, using the row sums and $a_i>0$ for $i<N$, gives $R_i(1)=1$ for $0\leq i\leq N$.
The terminal relation is $(x-b_N)R_N-c_NR_{N-1}=0$ on the orthogonality support. Linearization means equality of functions on this support, rather than an unrestricted polynomial identity when the degree sum exceeds $N$.

Write the tridiagonal recurrence matrix of such an array with the row convention
\[
J_{i,i+1}=a_i,\qquad J_{i,i}=b_i,\qquad J_{i,i-1}=c_i,
\]
with all other entries zero and indices outside $0,\ldots,N$ omitted.

\begin{lemma}[Finite recurrence realization and triangle support]\label{lem:finite}
Every recurrence satisfying \eqref{eq:basic} has a positive orthogonality measure on exactly $N+1$ real points. If $J$ is its tridiagonal recurrence matrix, then its finite product coefficients satisfy
\begin{equation}\label{eq:matrix-product}
g(m,n;k)=[R_m(J)]_{n,k}.
\end{equation}
They vanish unless $|m-n|\leq k\leq\min(m+n,N)$, regardless of their signs inside this interval.
\end{lemma}
\begin{proof}
Choose positive $\pi_i$ with $\pi_0=1$ and $\pi_{i+1}=\pi_i a_i/c_{i+1}$. For $D=\operatorname{diag}(\sqrt{\pi_i})$, $H=DJD^{-1}$ is real symmetric. Induction in \eqref{eq:recurrence} gives the row-vector identity $e_0^\mathsf T R_i(J)=e_i^\mathsf T$, hence $e_0^\mathsf T R_i(H)=\pi_i^{-1/2}e_i^\mathsf T$. In particular $e_0$ is cyclic for $H$. In an orthogonal eigenspace decomposition of $H$, the cyclic span of $e_0$ has dimension at most the number of distinct eigenvalues; here it is all of $\mathbb R^{N+1}$, so $H$ has $N+1$ distinct eigenvalues and the projection of $e_0$ on each eigenspace is nonzero. The spectral measure $\mu$ of $H$ at $e_0$ therefore has positive mass at each of $N+1$ real points, and
\[
\int R_iR_j\,d\mu=e_0^\mathsf TR_i(H)R_j(H)e_0
=\boldsymbol1_{i=j}\pi_i^{-1}.
\]
The terminal relation holds on this support: its polynomial $Q(x)=(x-b_N)R_N-c_NR_{N-1}$ satisfies $e_0^\mathsf TQ(J)=0$. Since $e_0$ is cyclic and $Q(J)$ commutes with $J$, $Q(J)=0$.
Multiplying a basis polynomial by $x$ acts by the corresponding row of $J$, including the terminal row. Induction on powers of $x$ therefore gives \eqref{eq:matrix-product}. Orthogonality also gives
\[
g(m,n;k)=\frac{\int R_mR_nR_k\,d\mu}{\int R_k^2\,d\mu}.
\]
If one of $m,n,k$ exceeds the sum of the other two, its polynomial is orthogonal to their product, whose degree is smaller. This proves the asserted support.
\end{proof}

We say that the array has \emph{nonnegative linearization}, or \emph{property (P)}, if $g(m,n;k)\geq0$ for all $0\leq m,n,k\leq N$. The coefficients $g(m,n;\cdot)$ then sum to one, since $R_i(1)=1$ and $1$ lies in the support, because $J\boldsymbol1=\boldsymbol1$.


For $q\geq1$ and $0<r\leq1$ let $G_{N,q,r}$ be the array
\begin{equation}\label{eq:parameters}
a_i=\frac{q^i[N-i]_q}{[N]_q},\qquad
b_i=(1-r)\frac{[i]_q}{[N]_q},\qquad
c_i=r\frac{[i]_q}{[N]_q},
\end{equation}
with $[j]_q$ as above; it satisfies \eqref{eq:basic}, since $[N]_q-[i]_q=q^i[N-i]_q$. Its profile $d_i=[i]_q/[N]_q$ has geometric increments with ratio $q$, and $b_i/c_i=(1-r)/r$ is constant. The Krawtchouk array is $K_{N,r}=G_{N,1,r}$, with $a_i=(N-i)/N$, $b_i=(1-r)i/N$ and $c_i=ri/N$.

\section{Log-concave profiles}\label{sec:lc}

For an array satisfying \eqref{eq:basic} with profile $d_i=b_i+c_i$ write
\[
\nu_k=\sum_{i=0}^k(-1)^{k-i}\binom ki d_i\quad(1\leq k\leq N),\qquad d_n=\sum_{k=1}^n\binom nk\nu_k\quad(0\leq n\leq N);
\]
below we put $\nu_0=0$ and $\nu_k=0$ for $k>N$.
We say that $(\nu_k)_{1\leq k\leq N}$ is \emph{log-concave without internal zeros} if all $\nu_k\geq0$, the indices with $\nu_k>0$ form an interval $\{1,\ldots,K\}$, and $\nu_k^2\geq\nu_{k-1}\nu_{k+1}$ for $2\leq k\leq N-1$. Note $\nu_1=d_1>0$.

\begin{theorem}[Log-concave Newton criterion]\label{thm:lc}
Let $0=d_0<d_1<\cdots<d_N=1$ have Newton coefficients that are log-concave without internal zeros, let $0<r\leq1$, and put $a_i=1-d_i$, $b_i=(1-r)d_i$, $c_i=rd_i$. Then the system has nonnegative linearization.
\end{theorem}

The hypotheses are inequalities in the profile. They hold for the Krawtchouk arrays ($\nu=(1/N,0,\ldots,0)$) and for all arrays $G_{N,q,r}$ ($\nu_k$ proportional to $(q-1)^{k-1}$, log-concave with equality). The hypotheses also hold for every admissible convex quadratic profile: $\nu_1>0$, $\nu_2\geq0$ when $N\geq2$, and $\nu_k=0$ for $3\leq k\leq N$. The hypotheses exclude the known counterexamples with constant ratio: the convex counterexamples recorded in~\cite{P3} have Newton coefficients of both signs, and for the combinations $(1-\lambda)K_{N,r}+\lambda G_{N,q,r}$ of Section~\ref{sec:limits}, with $q>1$, $N\geq3$ and $0<\lambda<1$, one has $\nu_1=(1-\lambda)/N+\lambda/[N]_q$ and $\nu_k=\lambda(q-1)^{k-1}/[N]_q$ for $2\leq k\leq N$, so $\nu_2^2-\nu_1\nu_3=-\lambda(1-\lambda)(q-1)^2/(N[N]_q)<0$.

\paragraph{A rescaled system and a kernel.}
Fix $r$ and put $t=(1-r)/\sqrt r\geq0$. If $p_i$ is defined by $p_0=1$ and $yp_i=a_ip_{i+1}+td_ip_i+d_ip_{i-1}$ $(0\leq i<N)$, then $R_i(x)=r^{i/2}p_i(x/\sqrt r)$, so the coefficients $H_{mn}^k$ in $p_mp_n=\sum_kH_{mn}^kp_k$ on the support satisfy $g(m,n;k)=r^{(m+n-k)/2}H_{mn}^k$. It suffices to show $H_{mn}^k\geq0$. Put
\[
e_i=\prod_{j=0}^{i-1}\frac{a_j}{d_{j+1}},\qquad E_i=e_ip_i.
\]
Then $yE_i=d_{i+1}E_{i+1}+td_iE_i+a_{i-1}E_{i-1}$ on the support for $0\leq i\leq N$, with the terms $d_{N+1}E_{N+1}$ and $a_{-1}E_{-1}$ omitted; at $i=N$ this is the terminal relation. Write $E_mE_n=\sum_kL_{mn}^kE_k$, so that $H_{mn}^k=L_{mn}^ke_k/(e_me_n)$, and
\[
K_k(z,w)=\sum_{m,n=0}^NL_{mn}^kz^mw^n.
\]
Let $\theta_z=z\,\partial_z$ and let $d(\theta)=\sum_k\nu_k\binom\theta k$ act on polynomials by $d(\theta_z)z^i=d_iz^i$. For $F(z)=\sum_{i=0}^NE_iz^i$ the recurrence reads $yF(z)=\mathcal L_zF(z)$ with
\[
\mathcal L_z=(z^{-1}+t-z)\,d(\theta_z)+z,
\]
because $d_0=0$ removes the term $z^{-1}E_0$ and $1-d_N=0$ removes the term $z^{N+1}E_N$. Applying this to $F(z)F(w)=\sum_kK_k(z,w)E_k$ in the two variables separately and comparing coefficients of the basis $E_k$ gives $(\mathcal L_z-\mathcal L_w)K_k=0$. Since $K_k$ is symmetric, $K_k=\sum_{\ell,a\geq0}\omega_{\ell,a}S^aZ^\ell$ with $S=z+w$, $Z=zw$, and $K_k(z,0)=z^k$ gives $\omega_{0,a}=\mathbf1_{a=k}$. As $K_k$ has degree at most $N$ in $z$, comparing the coefficients of $z^{\ell+a}w^\ell$ by induction on $\ell$ gives $\omega_{\ell,a}=0$ whenever $\ell+a>N$.

\paragraph{The kernel recurrence.}
Put $f_n(\ell)=\sum_{h\geq0}\binom\ell h\nu_{n+h}$ with $\nu_0=0$ and $\nu_k=0$ for $k>N$, so $f_0(\ell)=d_\ell$, and
\[
W_{p,q}(\ell,a)=\sum_{n\geq p}\binom{n-p+q}q\binom an f_n(\ell)\geq0,
\]
and $A_j=W_{2j+1,j}$, $B_j=W_{2j+2,j}$, $C_j=W_{2j+2,j+1}$.

\begin{lemma}\label{lem:kernel-rec}
For all $\ell,a\geq0$,
\begin{equation}\label{eq:kernel-rec}
\begin{aligned}
d_{\ell+1}\omega_{\ell+1,a}={}&a_{\ell+a}\,\omega_{\ell,a}+\sum_{j\geq0}(-1)^j\Big(t\,A_j(\ell-j,a+2j+1)\,\omega_{\ell-j,a+2j+1}\\
&+B_j(\ell-j,a+2j+2)\,\omega_{\ell-j,a+2j+2}\\
&+C_j(\ell-j-1,a+2j+2)\,\omega_{\ell-j-1,a+2j+2}\Big),
\end{aligned}
\end{equation}
where $a_i=1-d_i$ is extended by the polynomial $d_i$ for $i>N$ and terms with a negative index are omitted.
\end{lemma}
\begin{proof}
Since $S^aZ^\ell=\sum_i\binom ai z^{\ell+i}w^{\ell+a-i}$, Vandermonde's identity $\binom{\ell+i}k=\sum_{p+q=k}\binom\ell q\binom ip$ and $\binom ai\binom ip=\binom ap\binom{a-p}{i-p}$ give
\[
d(\theta_z)\,S^aZ^\ell=\sum_{p,q\geq0}\nu_{p+q}\binom\ell q\binom ap\,z^pS^{a-p}Z^\ell .
\]
Hence $(\mathcal L_z-\mathcal L_w)S^aZ^\ell$ equals $(z-w)S^aZ^\ell$ plus $\sum_{p,q}\nu_{p+q}\binom\ell q\binom ap S^{a-p}Z^\ell$ times $(z^{p-1}-w^{p-1})+t(z^p-w^p)-(z^{p+1}-w^{p+1})$. Divide by $z-w$, using $\chi_n=(z^{n+1}-w^{n+1})/(z-w)=\sum_{0\leq j\leq\lfloor n/2\rfloor}(-1)^j\binom{n-j}jS^{n-2j}Z^j$ for $n\geq0$, $\chi_{-1}=0$ and $(z^{-1}-w^{-1})/(z-w)=-Z^{-1}$. Since $\sum_q\binom\ell q\nu_{p+q}=f_p(\ell)$, the quotient is
\[
\begin{aligned}
&-d_\ell S^aZ^{\ell-1}+a_{\ell+a}S^aZ^\ell\\
&+\sum_{j\geq0}(-1)^j\Big(tA_j(\ell,a)S^{a-2j-1}Z^{\ell+j}+B_j(\ell,a)S^{a-2j-2}Z^{\ell+j}+C_j(\ell,a)S^{a-2j-2}Z^{\ell+j+1}\Big).
\end{aligned}
\]
Here $-d_\ell S^aZ^{\ell-1}$ comes from $p=0$; the term $a_{\ell+a}S^aZ^\ell$ combines $S^aZ^\ell$ with the parts $j=0$ of $-\chi_p$, which sum to $-d_{\ell+a}S^aZ^\ell$ by Vandermonde; and the three families from $\chi_{p-2}$, $t\chi_{p-1}$ and $-(\chi_p-S^p)$, whose coefficients $\binom{p-2-j}j$, $\binom{p-1-j}j$ and $\binom{p-1-j}{j+1}$ are those of $B_j$, $A_j$ and $C_j$. Summing over $\ell,a$ with weights $\omega_{\ell,a}$ and taking the coefficient of $S^aZ^\ell$ gives \eqref{eq:kernel-rec}, because the monomials $S^aZ^\ell$ are linearly independent.
\end{proof}

\paragraph{The log-concavity inequality.}
\begin{lemma}\label{lem:lc-ineq}
Under the hypotheses of Theorem~\ref{thm:lc}, for $i,l\geq1$ and $\ell\geq0$,
\[
f_i(\ell)\,f_l(\ell+1)\geq d_{\ell+1}\,f_{i+l}(\ell).
\]
\end{lemma}
\begin{proof}
For a nonnegative finitely supported sequence $u$ that is log-concave without internal zeros, $v_i=u_i+u_{i+1}$ has the same property, since
$v_i^2-v_{i-1}v_{i+1}=(u_i^2-u_{i-1}u_{i+1})+(u_{i+1}^2-u_iu_{i+2})+(u_iu_{i+1}-u_{i-1}u_{i+2})$
and the last term is nonnegative by log-concavity. The sequence $(0,\nu_1,\nu_2,\ldots)$ has this property, and $f_i(\ell)$ is obtained from it by applying $u\mapsto v$ $\ell$ times. For such a sequence $g$, balanced products dominate: $g_ig_l\geq g_0g_{i+l}$ and $g_ig_{l+1}\geq g_1g_{i+l}$ for $i,l\geq1$. Adding these with $g=f(\ell)$ and using $g_l+g_{l+1}=f_l(\ell+1)$ and $g_0+g_1=f_0(\ell+1)=d_{\ell+1}$ gives the claim.
\end{proof}

\begin{lemma}\label{lem:detours}
Under the hypotheses of Theorem~\ref{thm:lc}, for all $j\geq1$ and $\ell,a\geq0$,
\[
\begin{aligned}
A_0(\ell,a)B_{j-1}(\ell+1,a-1)+B_0(\ell,a)A_{j-1}(\ell+1,a-2)&\geq4d_{\ell+1}A_j(\ell,a),\\
B_0(\ell,a)B_{j-1}(\ell+1,a-2)&\geq4d_{\ell+1}B_j(\ell,a),\\
B_0(\ell,a)C_{j-1}(\ell+1,a-2)&\geq4d_{\ell+1}C_j(\ell,a),
\end{aligned}
\]
with $W_{p,q}(\ell,a)=0$ for $a<0$.
\end{lemma}
\begin{proof}
Use $A_0(\ell,a)\geq af_1(\ell)$ and $B_0(\ell,a)\geq\binom a2f_2(\ell)$, expand the second factor, apply Lemma~\ref{lem:lc-ineq} to $f_1(\ell)f_n(\ell+1)$ and $f_2(\ell)f_n(\ell+1)$, and use $a\binom{a-1}n=(n+1)\binom a{n+1}$ and $\binom a2\binom{a-2}n=\binom{n+2}2\binom a{n+2}$. After division by $d_{\ell+1}$ the three left sides are bounded below by $\sum_m\gamma_m\binom amf_m(\ell)$ with
\[
\gamma_m=\binom{m+1}2\binom{m-j-2}{j-1},\qquad\binom m2\binom{m-j-3}{j-1},\qquad\binom m2\binom{m-j-2}j,
\]
while the right sides are $4\sum_m\gamma'_m\binom amf_m(\ell)$ with $\gamma'_m=\binom{m-j-1}j$, $\binom{m-j-2}j$ and $\binom{m-j-1}{j+1}$. Where $\gamma'_m>0$ the quotients are
\[
\frac{jm(m+1)}{2(m-j-1)},\qquad\frac{jm(m-1)}{2(m-j-2)},\qquad\frac{(j+1)m(m-1)}{2(m-j-1)},
\]
on $m\geq2j+1$, $m\geq2j+2$ and $m\geq2j+2$ respectively. They increase with $j$ for fixed $m$; at $j=1$ their minima over integers are $5$, $5$ and $6$. All quotients exceed $4$, and $f_m(\ell)\geq0$.
\end{proof}

\begin{proof}[Proof of Theorem~\ref{thm:lc}]
Let $\Omega=\{(\ell,a):\ell,a\geq0,\ \ell+a\leq N\}$, on which $\omega$ is supported. Read \eqref{eq:kernel-rec} as moves into the node $(\ell+1,a)$, each divided by $d_{\ell+1}>0$: the straight move from $(\ell,a)$ with weight $a_{\ell+a}\geq0$ on $\Omega$, and for each $j$ the three moves from $(\ell-j,a+2j+1)$, $(\ell-j,a+2j+2)$, $(\ell-j-1,a+2j+2)$ with weights $tA_j$, $B_j$, $C_j$. Every move raises $\ell$. Let $\mathcal U$ be the straight moves, $\mathcal P$ the moves with $j$ even and $\mathcal Q$ the moves with $j$ odd, all with nonnegative weights, so that $\omega=\beta+(\mathcal U+\mathcal P-\mathcal Q)\omega$ with $\beta=\omega|_{\ell=0}$.

We claim $\mathcal P^2\geq4\mathcal Q$ entrywise. A move of $\mathcal Q$ with index $j$ from $x=(\ell,a)$ is matched by the two-step paths in $\mathcal P^2$ that start with $A_0$ or $B_0$ from $x$ (index $0$, landing at $(\ell+1,a-1)$ or $(\ell+1,a-2)$, which lie in $\Omega$) and continue with the move of the same type and index $j-1$ (even), as in Lemma~\ref{lem:detours}: $A_j$ by $A_0$ then $B_{j-1}$ and $B_0$ then $A_{j-1}$; $B_j$ by $B_0$ then $B_{j-1}$; $C_j$ by $B_0$ then $C_{j-1}$. These paths end at the target of the matched move, carry the same power of $t$, and have the extra denominator $d_{\ell+1}$ of the intermediate node. Distinct moves of $\mathcal Q$ have distinct displacements, so their paths are distinct, and Lemma~\ref{lem:detours} gives the claim.

Now show by induction on the level that $\omega\geq0$ and $\omega\geq\frac12\mathcal P\omega$. Both hold at $\ell=0$. If they hold below a node $y$ at level $\ell+1\geq1$, then, all weights being nonnegative, $(\mathcal P\omega)(y)\geq\frac12(\mathcal P^2\omega)(y)\geq2(\mathcal Q\omega)(y)$, and
\[
\omega(y)=(\mathcal U\omega)(y)+(\mathcal P\omega)(y)-(\mathcal Q\omega)(y)\geq\tfrac12(\mathcal P\omega)(y)\geq0.
\]
Hence all $\omega_{\ell,a}\geq0$, so $K_k$ has nonnegative coefficients in $z,w$, so $L_{mn}^k\geq0$, $H_{mn}^k\geq0$ and $g(m,n;k)\geq0$.
\end{proof}

\begin{remark}\label{rem:bernstein}
The argument is valid coefficientwise in $t$. With independent parameters in $yp_i=a_ip_{i+1}+ud_ip_i+vd_ip_{i-1}$, weighted homogeneity ($\deg y=\deg u=1$, $\deg v=2$) then shows that each linearization coefficient is a polynomial in $(u,v)$ with nonnegative coefficients: for $D=m+n-k\geq0$ it has the form $\sum_{0\leq j\leq\lfloor D/2\rfloor}h_ju^{D-2j}v^j$ with $h_j\geq0$, and it vanishes for $D<0$. Consequently every $g(m,n;k)$ of the array of Theorem~\ref{thm:lc}, as a function of $r$, equals $\sum_jh_jr^j(1-r)^{D-2j}$; multiplying the $j$-th term by $(r+(1-r))^j=1$ shows that its coefficients in the Bernstein basis of degree $D$ are $\binom Ds^{-1}\sum_{j\leq s\leq2j}h_j\binom j{s-j}\geq0$.
\end{remark}

\begin{corollary}[Full positive Krawtchouk range]\label{cor:kraw}
For every $N\geq1$ and $0<p_0\leq1/2$, Krawtchouk polynomials orthogonal for binomial weights $\binom Nj p_0^{N-j}(1-p_0)^j$ have nonnegative linearization when normalized at $j=0$.
\end{corollary}
\begin{proof}
To fix the convention without relying on a named recurrence, put $r=p_0/(1-p_0)\in(0,1]$, and let $Y_1,\ldots,Y_N$ be independent variables with values $1,-r$ of probabilities $p_0,1-p_0$. Their mean is zero and $Y_i^2=(1-r)Y_i+r$. The normalized elementary symmetric sums
\[
K_i(j)=\binom Ni^{-1}\sum_{|S|=i}\prod_{l\in S}Y_l
\]
depend only on the number $j$ of negative coordinates. They are orthogonal, because the expectation of two distinct squarefree monomials is zero: some independent mean-zero variable occurs once. Their weights in $j$ are exactly the stated binomial weights. With $x=N^{-1}\sum_lY_l=1-j/(N(1-p_0))$, multiplication by $x$ gives
\[
xK_i=\frac{N-i}{N}K_{i+1}+(1-r)\frac iN K_i+r\frac iN K_{i-1}.
\]
Indeed, in multiplying the elementary sum by $\sum_lY_l$, indices outside a subset give $(i+1)$ copies of each subset of size $i+1$; the linear part of $Y_l^2$ gives $(1-r)i$ copies of each subset of size $i$; and the constant part gives $r(N-i+1)$ copies of each subset of size $i-1$. Division by the binomial normalizers gives the displayed recurrence, including $i=N$. Its leading coefficient in $x$ is positive, its degree is $i$, and its value at $j=0$ is one. These properties identify the normalized orthogonal polynomials for the binomial measure. This is the array $K_{N,r}$; its profile $d_i=i/N$ has Newton coefficients $(1/N,0,\ldots,0)$, so Theorem~\ref{thm:lc} applies.
The range is sharp in this endpoint normalization: if $1/2<p_0<1$, then
$r=p_0/(1-p_0)>1$ and $g(1,1;1)=b_1=(1-r)/N<0$.
For $N=1$, this coefficient follows from the terminal relation.
\end{proof}

\begin{corollary}[Geometric arrays]\label{cor:additional}
For every $N\geq2$, $q>1$, and $0<r\leq1$, the array $G_{N,q,r}$ of \eqref{eq:parameters} has nonnegative linearization. In the normalization $R_0=1$, $R_1=x$, it is outside the Krawtchouk recurrence class.
\end{corollary}
\begin{proof}
Its increments are $d_i-d_{i-1}=q^{i-1}/[N]_q$, so its Newton coefficients $\nu_k=(q-1)^{k-1}/[N]_q$ are log-concave without internal zeros and Theorem~\ref{thm:lc} applies. For a Krawtchouk recurrence in this normalization the increments of $b_i+c_i$ are constant. Here their ratio is $q>1$ and $N\geq2$, so they are not constant. An affine change of the support transporting the same endpoint-normalized basis leaves the multiplication coefficients in the basis unchanged; it cannot remove this distinction.
\end{proof}

The arrays $G_{N,q,r}$ are the normalized recurrences of Koelink's construction~\cite[Theorems~3.1(ii) and~7.1]{koelink} with $r=1/p$; their polynomial interpretation is the dual $q$-Krawtchouk interpretation described after his Theorem~7.8. A proof of Corollary~\ref{cor:additional} by a finite algebra with positive multiplication, which needs only the geometric increments, is given in~\cite{P3}.

\section{Szwarc's quadratic system}\label{sec:szwarc}

Problem~2 of~\cite[\S5]{szwarc} does not stop at the Krawtchouk recurrence. With his parameter written as $p_0$, Szwarc introduces the symmetric recurrence of length $2N+1$
\begin{equation}\label{eq:szwarc-q}
yQ_{2n}=(N-n)Q_{2n+1}+nQ_{2n-1},\qquad
yQ_{2n-1}=(1-p_0)Q_{2n}+p_0Q_{2n-2},
\end{equation}
records that $Q_{2n}(y)$ is the Krawtchouk polynomial $K_n(y^2)$ of his recurrence, states that the $Q_n$ have nonnegative linearization for $p_0\leq1/2$, and asks for a general criterion that would handle this case. His Theorem~3.3 for infinite symmetric systems does not apply. In the normalization \eqref{eq:szwarc-q} its monotonicity hypotheses hold, and its hypothesis $\alpha_n\leq\gamma_n$ holds at odd indices exactly when $p_0\leq1/2$; at an even index $2n$ it reads $n\leq N-n$ and fails in the upper half. Some failure of this kind is forced in every finite system, whose terminal upward coefficient vanishes. We derive the positivity of \eqref{eq:szwarc-q} from Theorem~\ref{thm:lc} through a quadratic decomposition of symmetric finite systems.

Put $M=2N$, $R_{2n}(t)=Q_{2n}(\sqrt Nt)$ and $R_{2n-1}(t)=\sqrt N\,Q_{2n-1}(\sqrt Nt)$. Substitution in \eqref{eq:szwarc-q} shows that the $R_k$ satisfy \eqref{eq:recurrence}, with $N$ replaced by $M$, for the array
\begin{equation}\label{eq:qarray}
b_k=0,\qquad
(a_{2n},c_{2n})=\Big(\frac{N-n}{N},\frac nN\Big),\qquad
(a_{2n-1},c_{2n-1})=(1-p_0,\,p_0).
\end{equation}
It satisfies \eqref{eq:basic} with $M$ in place of $N$ whenever $0<p_0<1$. The two normalizations differ by positive factors and a positive dilation, so they have the same product coefficients up to positive factors.

Szwarc uses the quadratic transformation~\cite[\S3.5]{szwarc} to pass from a symmetric system to its even part. For finite systems the converse direction also works: positivity of a symmetric system follows from positivity of two nonsymmetric finite systems, both read off from its coefficients. Any criterion for nonsymmetric finite systems therefore becomes a criterion for symmetric ones.

\begin{theorem}[Quadratic decomposition]\label{thm:decomp}
Let $N\geq2$, and let $a_k,c_k$ $(0\leq k\leq2N)$ with $b\equiv0$ satisfy \eqref{eq:basic} with $2N$ in place of $N$. Put $a_k=c_k=0$ for $k\notin[0,2N]$, and define the even component
\[
\hat a_n=a_{2n}a_{2n+1},\quad \hat b_n=c_{2n}a_{2n-1}+a_{2n}c_{2n+1},\quad \hat c_n=c_{2n}c_{2n-1}\qquad(0\leq n\leq N)
\]
and the odd component
\[
\check a_n=a_{2n+1}a_{2n+2},\quad \check b_n=a_{2n+1}c_{2n+2}+c_{2n+1}a_{2n},\quad \check c_n=c_{2n+1}c_{2n}\qquad(0\leq n<N).
\]
After the affine normalization $(\hat a_n,\hat b_n,\hat c_n)\mapsto(\hat a_n,\hat b_n-\hat b_0,\hat c_n)/\hat a_0$, and likewise for the odd component, these arrays satisfy \eqref{eq:basic} on the indices $0,\ldots,N$ and $0,\ldots,N-1$. If both have nonnegative linearization, so does the symmetric array. More precisely, with $\hat g$ and $\check g$ the product coefficients of the components and terms with indices out of range omitted,
\[
g(2a,2b;2c)=\hat g(a,b;c),\qquad
g(2a,2b+1;2c+1)=a_{2a}\,\check g(a,b;c)+c_{2a}\,\check g(a-1,b;c),
\]
the coefficients with $m+n+k$ odd vanish, and $g(2a+1,2b+1;2c)\int R_{2c}^2\,d\mu=g(2c,2a+1;2b+1)\int R_{2b+1}^2\,d\mu$. Conversely, nonnegative linearization of the symmetric array implies that of its even component.
\end{theorem}
\begin{proof}
Since $b\equiv0$, the matrix $J$ is conjugate to $-J$ by $\operatorname{diag}((-1)^i)$, which fixes $e_0$; hence the measure of Lemma~\ref{lem:finite} is symmetric and $R_k(-t)=(-1)^kR_k(t)$. By symmetry $R_{2n}(t)=P_n(t^2)$ and $R_{2n+1}(t)=tL_n(t^2)$ for polynomials $P_n$, $L_n$ of degree $n$. Applying \eqref{eq:recurrence} twice gives, with $x=t^2$,
\[
xP_n=\hat a_nP_{n+1}+\hat b_nP_n+\hat c_nP_{n-1},\qquad
xL_n=\check a_nL_{n+1}+\check b_nL_n+\check c_nL_{n-1}.
\]
The row sums of the components equal $a_{2n}+c_{2n}=1$ and $a_{2n+1}+c_{2n+1}=1$; the coefficients $\hat a_n,\check a_n$ are positive below the last index and vanish there, and $\hat c_n,\check c_n>0$ for $n\geq1$. Hence the normalized arrays, which describe the same polynomials in the variable $(x-\hat b_0)/\hat a_0$ (respectively $(x-\check b_0)/\check a_0$), satisfy \eqref{eq:basic}. The terminal relation $tR_{2N}=c_{2N}R_{2N-1}$ on the support $\Sigma$ of the symmetric system, applied once and twice, gives the terminal relations of $P$ on $\{t^2:t\in\Sigma\}$ and of $L$ on $\{t^2:t\in\Sigma,\ t\neq0\}$. These sets have $N+1$ and $N$ points, because $\Sigma$ is symmetric with $2N+1$ points and contains $0$. By Lemma~\ref{lem:finite} they are the supports of the two components.

The identity $tR_{2n}=a_{2n}R_{2n+1}+c_{2n}R_{2n-1}$ gives $P_n=a_{2n}L_n+c_{2n}L_{n-1}$ on the support of $L$; for $n=N$ it is the terminal relation, with $a_{2N}=0$. Hence, on $\Sigma\setminus\{0\}$,
\[
R_{2a}R_{2b+1}=t\big(a_{2a}L_aL_b+c_{2a}L_{a-1}L_b\big)=\sum_c\big(a_{2a}\check g(a,b;c)+c_{2a}\check g(a-1,b;c)\big)R_{2c+1},
\]
and both sides vanish at $t=0$. Similarly $R_{2a}R_{2b}=(P_aP_b)(t^2)=\sum_c\hat g(a,b;c)R_{2c}$ on $\Sigma$. Uniqueness of expansions on $\Sigma$ gives the displayed formulas. The vanishing for $m+n+k$ odd and the relation for two odd indices follow from the symmetry of the measure and from $g(m,n;k)\int R_k^2\,d\mu=\int R_mR_nR_k\,d\mu$. The converse holds because $\hat g(a,b;c)=g(2a,2b;2c)$.
\end{proof}

Combined with Theorem~\ref{thm:lc}, the decomposition gives a parity-split criterion of the same shape as Szwarc's Theorem~3.3: separate conditions on the even-indexed and the odd-indexed coefficients.

\begin{corollary}[Parity-split criterion]\label{cor:parity}
Let $N\geq1$ and let a symmetric array $(b\equiv0)$ on the indices $0,\ldots,2N$ satisfy \eqref{eq:basic}. Suppose that
\begin{enumerate}
\item the odd downward coefficients are constant, $c_{2n+1}=p_0$, with $0<p_0\leq1/2$;
\item the even downward coefficients $d_n=c_{2n}$ have positive increments $d_n-d_{n-1}$, $1\leq n\leq N$, forming a geometric progression with ratio $q\geq1$.
\end{enumerate}
Then the array has nonnegative linearization. With $r=p_0/(1-p_0)$ and $N\geq2$, its normalized even and odd components are \eqref{eq:parameters} with parameters $(N,q,r)$ and $(N-1,q,r/q)$, respectively. For $N=1$ the odd component consists only of $L_0=1$ and the array is \eqref{eq:qarray}. If $1/2<p_0<1$ in (i) and (ii) holds, then $g(2,2;2)<0$.
\end{corollary}
\begin{proof}
For $N=1$ the array is $a_0=1$, $c_1=p_0$, $a_1=1-p_0$, $c_2=1$, with support $\{-1,0,1\}$ and $R_1=t$, $R_2=(t^2-p_0)/(1-p_0)$. On the support $R_1^2=(1-p_0)R_2+p_0$, $R_1R_2=R_1$ and $R_2^2=\frac{p_0}{1-p_0}+\frac{1-2p_0}{1-p_0}R_2$, as evaluation at $t=0,\pm1$ shows; all coefficients are nonnegative for $p_0\leq1/2$, and $g(2,2;2)<0$ for $p_0>1/2$. Let $N\geq2$. The even component has $\hat a_n=(1-d_n)(1-p_0)$, $\hat b_n=d_n(1-p_0)+(1-d_n)p_0$ and $\hat c_n=d_np_0$, so $\hat a_0=1-p_0$, $\hat b_0=p_0$, and its normalization is $(1-d_n,\,(1-r)d_n,\,rd_n)$. Its profile $d_n$ has geometric increments, so its Newton coefficients are proportional to $(q-1)^{k-1}$ and log-concave without internal zeros, and $0<r\leq1$. For the odd component, $\check a_n=(1-p_0)(1-d_{n+1})$, $\check b_n=(1-p_0)d_{n+1}+p_0(1-d_n)$ and $\check c_n=p_0d_n$. Geometric increments give $d_{n+1}-d_1=qd_n$, so the normalization is
\[
\Big(\frac{1-d_{n+1}}{1-d_1},\ \frac{(q-r)d_n}{1-d_1},\ \frac{rd_n}{1-d_1}\Big).
\]
Its profile $qd_n/(1-d_1)$ has geometric increments with ratio $q$, and its downward coefficients are $r/q\leq1$ times the profile. Theorem~\ref{thm:lc} applies to both components, and Theorem~\ref{thm:decomp} concludes. Comparison with \eqref{eq:parameters}, using $d_n=[n]_q/[N]_q$ and $1-d_1=q[N-1]_q/[N]_q$, identifies the parameters. For $p_0>1/2$, $g(2,2;2)=\hat g(1,1;1)$ is the normalized diagonal coefficient $(1-r)d_1<0$ (for $N=1$ see above).
\end{proof}

\begin{theorem}[Szwarc's quadratic system]\label{thm:szwarc}
Let $N\geq1$ and $0<p_0\leq1/2$. The array \eqref{eq:qarray} has a symmetric positive orthogonality measure on $2N+1$ points, and for all $0\leq m,n\leq 2N$
\[
R_mR_n=\sum_{k=0}^{2N}g(m,n;k)R_k,\qquad g(m,n;k)\geq0,\qquad\sum_k g(m,n;k)=1
\]
on that support. If $1/2<p_0<1$, then $g(2,2;2)<0$.
\end{theorem}
\begin{proof}
Lemma~\ref{lem:finite} gives a positive measure on $2N+1$ points, and it is symmetric because $b\equiv0$ (proof of Theorem~\ref{thm:decomp}). The array \eqref{eq:qarray} satisfies the hypotheses of Corollary~\ref{cor:parity} with $c_{2n+1}=p_0$ and $d_n=n/N$, that is, with $q=1$, which gives nonnegative linearization for $p_0\leq1/2$ and $g(2,2;2)<0$ for $1/2<p_0<1$. The coefficients sum to one because $R_k(1)=1$ and $1$ lies in the support, since $J\boldsymbol1=\boldsymbol1$.
\end{proof}

The even component of \eqref{eq:qarray} is the Krawtchouk array $K_{N,r}$ with $r=p_0/(1-p_0)$, so that $Q_{2n}(y)$ is a Krawtchouk polynomial in $y^2$, as Szwarc notes, and its odd component is $K_{N-1,r}$. A second proof of Theorem~\ref{thm:szwarc}, by a graded version of the algebra of~\cite{P3}, and explicit coefficient neighbourhoods of \eqref{eq:qarray} are given there.

\section{How far can such criteria go?}\label{sec:limits}

A criterion valid for all $N$ cannot leave a margin of fixed size around the Krawtchouk arrays. Let
\[
\operatorname{dist}(\widetilde a,\widetilde b,\widetilde c)=\max_{0\leq i\leq N}\max\Bigl(\Bigl|\widetilde a_i-\frac{N-i}N\Bigr|,\Bigl|\widetilde b_i-\frac{(1-r)i}N\Bigr|,\Bigl|\widetilde c_i-\frac{ri}N\Bigr|\Bigr)
\]
be the coefficient distance to $K_{N,r}$.

\begin{lemma}[Corner value]\label{lem:corner-value}
For $0<r\leq1$, the array $K_{N,r}$ has $g(N,N;N)=(1-r)^N$.
\end{lemma}
\begin{proof}
Use the model in the proof of Corollary~\ref{cor:kraw}, with $p_0=r/(1+r)$. There $R_N=K_N=Y_1\cdots Y_N$ and $K_k=\binom Nk^{-1}\sum_{|S|=k}Y_S$, where $Y_S=\prod_{l\in S}Y_l$. Since $Y_l^2=(1-r)Y_l+r$,
\[
R_N^2=\prod_{l=1}^N\bigl((1-r)Y_l+r\bigr)=\sum_S(1-r)^{|S|}r^{N-|S|}Y_S.
\]
Both sides of $R_N^2=\sum_kg(N,N;k)R_k$ depend only on the number $j$ of negative coordinates, and the identity holds on the support, that is, for every $j$; so it holds as an identity of functions of $(Y_1,\ldots,Y_N)$. The $2^N$ monomials $Y_S$ are orthogonal for the product measure, with $\mathbb E[Y_S^2]=r^{|S|}>0$, hence linearly independent. The monomial $Y_{\{1,\ldots,N\}}$ occurs only in $K_N$, with coefficient one, which gives $g(N,N;N)=(1-r)^N$.
\end{proof}

Thus the corner coefficient vanishes at $r=1$. The next proposition raises the top mass of the orthogonality measure.

\begin{proposition}[Raising the top mass]\label{prop:raise}
Let $N\geq1$ and $0<r<1$. The measure of $K_{N,r}$ has nodes $x_n=1-(1+r)n/N$ and masses $w_n=\binom Nnr^{N-n}/(1+r)^N$. For $\tau>0$ let $\mu^\tau$ be this measure with the mass at $x_0=1$ raised to $(1+\tau)w_0$. Then $\mu^\tau$ has
\[
g(N,N;N)=\frac{(1+\tau)^2(I-w_0)+w_0}{(1+\tau)(h-w_0)+w_0},\qquad I=\bigl(r(1-r)\bigr)^N,\quad h=r^N,
\]
which is negative exactly when $\tau>\tau_N^*=\bigl(1-(1-r^2)^N\bigr)^{-1/2}-1$. After the affine normalization $R_n(1)=1$, $R_1=x$, its recurrence coefficients differ from those of $K_{N,r}$ by at most $3\tau\max_nw_n$.
\end{proposition}
\begin{proof}
Write $\pi_n$ and $S_n=\sum_{j<n}\pi_j$ for $K_{N,r}$, and put $\varepsilon=\tau w_0$. The polynomial $(1+\varepsilon S_n)R_n-\varepsilon\sum_{j<n}\pi_jR_j$ has degree $n$ and equals $1$ at $x_0$. For $\deg p<n$, its integral against $p$ with respect to $\mu^\tau=\mu+\varepsilon\delta_{x_0}$ is $-\varepsilon p(x_0)+\varepsilon p(x_0)=0$, because $\sum_{j<n}\pi_jR_j$ reproduces $p(x_0)$. Hence it is the polynomial $R_n^\tau$ of $\mu^\tau$.

On the support, dual orthogonality gives $\sum_{j\leq N}\pi_jR_j(x)=\boldsymbol1_{x=x_0}/w_0$, and in particular $S_N+\pi_N=1/w_0$. Therefore $R_N^\tau=(1+\tau)R_N-\tau\boldsymbol1_{x_0}$ on the support. By Lemma~\ref{lem:corner-value} and $\pi_N=r^{-N}$ (Lemma~\ref{lem:finite}), $\int R_N^3\,d\mu=I$; also $\int R_N^2\,d\mu=h$. Hence
\begin{gather*}
\int(R_N^\tau)^3d\mu^\tau=(1+\tau)^3(I-w_0)+(1+\tau)w_0,\\
\int(R_N^\tau)^2d\mu^\tau=(1+\tau)^2(h-w_0)+(1+\tau)w_0,
\end{gather*}
and their ratio is the stated formula. Since $r>r/(1+r)>r(1-r)$, we have $h>w_0>I$. So the denominator is positive, and the numerator is negative exactly when $(1+\tau)^2>w_0/(w_0-I)=\bigl(1-(1-r^2)^N\bigr)^{-1}$.

For the coefficients, comparing leading coefficients and using detailed balance, $\pi_na_n=\pi_{n+1}c_{n+1}$, for both measures gives
\[
a_n^\tau=a_n\frac{1+\varepsilon S_n}{1+\varepsilon S_{n+1}},\qquad c_n^\tau=c_n\frac{1+\varepsilon S_{n+1}}{1+\varepsilon S_n},
\]
and the row sums stay $1$. Since $\varepsilon\pi_n=\tau w_n$, each of $a_n$, $b_n$, $c_n$ changes by at most $\tau w_n$. Put $A_n=1+\varepsilon S_n$ and $t_n=\varepsilon\pi_n=\tau w_n$. Since $b_0^\tau=\varepsilon/(1+\varepsilon)$, the affine coordinate is $y=(1+\varepsilon)x-\varepsilon$, and the normalized coefficients are $\widetilde a_n=(1+\varepsilon)a_n^\tau$, $\widetilde b_n=(1+\varepsilon)b_n^\tau-\varepsilon$, $\widetilde c_n=(1+\varepsilon)c_n^\tau$. For $n\geq1$ we have $A_n\geq1+\varepsilon$, and
\begin{align*}
\widetilde a_n-a_n&=\varepsilon a_n-(1+\varepsilon)a_nt_n/A_{n+1},\\
\widetilde c_n-c_n&=\varepsilon c_n+(1+\varepsilon)c_nt_n/A_n,\\
\widetilde b_n-b_n&=-\varepsilon(a_n+c_n)+(1+\varepsilon)t_n\bigl(a_n/A_{n+1}-c_n/A_n\bigr).
\end{align*}
Because $a_n,c_n\geq0$ and $a_n+c_n\leq1$, each absolute change is at most $\varepsilon+t_n\leq2\tau\max_jw_j$. Row zero is restored exactly. This proves, in particular, the stated bound $3\tau\max_jw_j$ for every $\tau>0$.
\end{proof}

\begin{corollary}[Exponentially small stability radius]\label{cor:radius-exp}
Let $0<r<1$, and let $\rho_N(r)$ be the supremum of the $\rho$ such that every array satisfying \eqref{eq:basic}, with $\widetilde b_i\geq0$ for all $i$ and with $\operatorname{dist}(\widetilde a,\widetilde b,\widetilde c)\leq\rho$, has property (P). If $(1-r^2)^N\leq\tfrac12$ and $3(1-r^2)^N\leq(1-r)/N$, then
\[
\rho_N(r)\leq3\tau_N^*\max_nw_n\leq3(1-r^2)^N.
\]
In particular $\rho_N(r)=O\bigl((1-r^2)^N\bigr)$ as $N\to\infty$, for every fixed $r\in(0,1)$.
\end{corollary}
\begin{proof}
Put $x=(1-r^2)^N$ and $W=\max_nw_n$. For $0<x\leq\tfrac12$,
\[
(1+x)^2(1-x)-1=x(1-x-x^2)>0,
\]
so $\tau_N^*<x$. Choose $\tau\in(\tau_N^*,x)$. Proposition~\ref{prop:raise} gives a negative corner coefficient and full-array distance at most $3\tau W<3x\leq(1-r)/N$. Thus $\widetilde b_n>0$ for $n\geq1$, while $\widetilde b_0=0$. The exact coefficient formulas and the positive affine scale give $\widetilde a_n>0$ for $n<N$, $\widetilde c_n>0$ for $n>0$, and all the basic boundary and row-sum equalities. Letting $\tau\downarrow\tau_N^*$ proves $\rho_N(r)\leq3\tau_N^*W\leq3x$. For each fixed $r\in(0,1)$, both smallness conditions hold for all sufficiently large $N$, since $N(1-r^2)^N\to0$.
\end{proof}

Thus, for every fixed $r\in(0,1)$, no set of arrays with property (P) contains coefficient balls of a positive radius independent of $N$ around all $K_{N,r}$. For each fixed $N$, in contrast, $\rho_N(r)>0$ for $0<r<1$: explicit perturbation inequalities give $\rho_N(r)\geq\frac12(\min(r,1-r)/(32N))^N$, and, for fixed $N\geq2$, $\rho_N(r)$ is of exact order $(1-r)^N$ as $r\to1$; at the symmetric endpoint $r=1$, criteria that admit polynomial deformation rules, or analytic rules of the kind specified there, must satisfy two first-order equalities. These results are proved in~\cite{P3}.

\paragraph{How sharp is log-concavity?}
In Theorem~\ref{thm:lc} the Krawtchouk arrays lie on the edge of the hypothesis on the profile, and the arrays $G_{N,q,r}$ lie on its boundary, with equality in every log-concavity inequality. Positivity of the Newton coefficients of the profile does not suffice. For $q>1$ and $0<\lambda<1$ the combination $(1-\lambda)K_{N,r}+\lambda G_{N,q,r}$ is a stochastic array with constant ratio $b_i/c_i$ and increments proportional to $1+t\,q^{i-N}$, where $t=\lambda N(q-1)q^{N-1}/((1-\lambda)(q^N-1))$; its Newton coefficients are positive, and for $N\geq3$ one has $\nu_2^2<\nu_1\nu_3$ (Section~\ref{sec:lc}). Exact computation (\path{convexity_probe.py}) gives $g(N,N;N)<0$ for this combination at $(r,N,q,t)=(4/5,10,10,1)$, $(3/5,20,10,2)$, $(1/2,20,10,1)$, $(2/5,60,10,2)$, $(1/3,100,10,2)$, $(3/10,200,2,5)$ and $(2/7,200,10,5)$. In particular, at these parameters no sufficient criterion whose admitted set is convex, such as a system of linear inequalities in the coefficients, can admit both $K_{N,r}$ and $G_{N,q,r}$.

Log-concavity is not necessary either. For $0<\mu<1$ and $q>1$ let $D(i)=(1-\mu)i+\mu[i]_q$, whose Newton coefficients are $\nu_1=1$ and $\nu_k=\mu(q-1)^{k-1}$ for $k\geq2$; thus $\nu_1\nu_3/\nu_2^2=1/\mu$, and the other log-concavity inequalities hold with equality. For $\mu=\frac1{10}$ and $q=\frac{11}{10}$, the constant-ratio arrays with profile $d_i=D(i)/D(N)$ have property (P) for every $r\in(0,1]$ and $N\in\{6,8,10,12\}$ (machine-verified as \path{Szkraw.lc_violation_holds}). Every coefficient $g(m,n;k)$ is a polynomial in $r$ with nonnegative Bernstein coefficients (\path{runs/sharp/lc_violation_cert.py}). For $\mu=\frac1{100}$ and the same $q$, already $N=6$ gives a negative coefficient $g(6,6;6)$ at $r=1$. At $r=1$, float64 screening with high-precision confirmation for $N\leq20$ and high-precision searches in this family for $20\leq N\leq80$, both over coarse grids of $\mu$ and $q$, found as smallest factor $1/\mu$ breaking (P) $10^{5/2}$ at $N=4$, $10^2$ for $6\leq N\leq30$, and $10^{3/2}$ at $N=40,60,80$ (\path{runs/sharp/r1_grid.py}, \path{runs/sharp/r1_large.py}, \path{runs/sharp/r1_trend.py}; commands and outputs in \path{runs/sharp/RESULTS.md}). These are the smallest failures found, not proven thresholds: the observed failures lay in narrow windows of $(\mu,q)$ on the tested grids. Determine the exact boundary between log-concavity and positivity of the Newton coefficients of the profile. In particular, with $V_N=\inf\{\mu^{-1}:0<\mu<1,\ q>1,\ \text{the array with profile }D/D(N)\text{ and }r=1\text{ fails (P)}\}$ ($\inf\varnothing=\infty$), decide whether $V_N\to1$ as $N\to\infty$.

\section{Open problems}\label{sec:open}

\paragraph{The exact boundary.} Determine the set of profiles for which every constant-ratio array has property (P), or the threshold $V_N$ of Section~\ref{sec:limits}. The examples above suggest that the failures of log-convex profiles are driven by an acceleration of the increments near the terminal index.

\paragraph{Other ratios.} Theorem~\ref{thm:lc} requires $b_i/c_i$ to be constant. The companion paper~\cite{P2} admits ratios $c_i/d_i$ that increase with the index: linearly or as a convex quadratic on every log-concave Newton profile, with Newton coefficients of the ratio that are log-concave in every degree, and for cubic ratios under a single inequality, in each case when $b$ is nondecreasing. Criteria for other nonconstant ratios are open.

\paragraph{Monotonicity-type criteria.} Szwarc's Theorems~3.2 and~3.3 in~\cite{szwarc} are open conditions of monotonicity type for infinite systems. By Theorem~\ref{thm:decomp}, any such criterion for nonsymmetric finite systems, valid including the terminal indices, gives one for symmetric finite systems. Inequality conditions of this type for all products except those whose pairwise index sums all exceed $N$ are given in~\cite{P3}; for the remaining, deep products, Theorem~\ref{thm:lc} and its extensions in~\cite{P2} give shape conditions that control them.

\section*{AI-use disclosure}

OpenAI Codex and Anthropic Claude were used to develop proof arguments, check source passages, construct exact-arithmetic diagnostics, draft and revise this manuscript, and prepare the Lean formalization. The kernel recurrence and the operator inequality in the proof of Theorem~\ref{thm:lc} arose in research consultations with ChatGPT Pro; the closed form of the recurrence for every degree and the exact checks were completed with Claude, which also developed Sections~\ref{sec:szwarc} and~\ref{sec:limits} and prepared the Lean formalization of all numbered results. The review of Theorem~\ref{thm:lc} was therefore carried out by the same system that proposed its key steps. The results were reviewed in fourteen rounds by ChatGPT Pro, as part of a longer manuscript (an AI review with independent exact checks, not human peer review); the first round failed, the corrections are recorded in \path{REVIEW.md}, and the final confirmation rounds returned PASS. The remark on the sharpness of log-concavity in Section~\ref{sec:limits} passed a further review round with minor fixes. The author set the research objective and retains responsibility for publication decisions and the final mathematical claims. No independent human verification is asserted here. ``Machine-verified'' refers only to the Lean theorems listed in the section on code and formalization, with the axioms stated there; unlisted remarks and computational observations are not claimed machine-verified.

\section*{Code and formalization}

The companion code~\cite{code} contains exact-arithmetic diagnostics for every result (\path{papers/DIAGNOSTICS.md}) and a Lean~4 development. Run \verb|make lean| to build it and \verb|make axioms| to build and audit the principal statements; the audit invokes \path{lean/CheckAxioms.lean}, which rejects dependencies outside \path{propext}, \path{Classical.choice} and \path{Quot.sound}. All declarations below have compiled proofs without added axioms or admitted goals. Notes on how each Lean statement matches the text are in \path{papers/FORMALIZATION.md}; this file and \path{papers/DIAGNOSTICS.md} are in the later revision~\cite{docs}.
\begin{center}
\begin{tabular}{@{}p{0.36\linewidth}p{0.59\linewidth}@{}}
Result & Lean declaration\\[3pt]
Lemma~\ref{lem:finite}: finite realization & \path{Szkraw.finite_realization}, \path{Szkraw.product_triangle_zero}\\[3pt]
Theorem~\ref{thm:lc}: log-concave Newton criterion & \path{Szkraw.logconcave_criterion}\\[3pt]
Remark~\ref{rem:bernstein} & \path{Szkraw.lc_twoParam}, \path{Szkraw.lc_bernstein}, \path{Szkraw.constantRatio_twoParam}\\[3pt]
Corollary~\ref{cor:kraw}: Krawtchouk range & \path{Szkraw.full_krawtchouk_binomial_range}, \path{Szkraw.binomial_normalized_at_zero}, \path{Szkraw.binomial_node_probability_convention}, \path{Szkraw.krawtchouk_range_sharp}\\[3pt]
Corollary~\ref{cor:additional}: geometric arrays & \path{Szkraw.additional_geometric_class}\\[3pt]
Theorem~\ref{thm:decomp}: quadratic decomposition & \path{Szkraw.decomp_positivity}, \path{Szkraw.decomp_converse}, \path{Szkraw.even_block}, \path{Szkraw.odd_block}, \path{Szkraw.evenComp_admissible}, \path{Szkraw.oddComp_admissible}\\[3pt]
Corollary~\ref{cor:parity}: parity-split criterion & \path{Szkraw.parity_split}, \path{Szkraw.parity_split_sharp}, \path{Szkraw.ps_even_params}, \path{Szkraw.ps_odd_params}\\[3pt]
Theorem~\ref{thm:szwarc}: Szwarc's quadratic system & \path{Szkraw.szwarc_full}, \path{Szkraw.szwarc_sharp}, \path{Szkraw.realization_symmetric}\\[3pt]
Lemma~\ref{lem:corner-value}: corner value & \path{Szkraw.corner_value}\\[3pt]
Proposition~\ref{prop:raise}: raising the top mass & \path{Szkraw.raise_measure}, \path{Szkraw.raise_corner}, \path{Szkraw.raise_sign}, \path{Szkraw.raise_bound}\\[3pt]
Corollary~\ref{cor:radius-exp} and its no-uniform-radius conclusion & \path{Szkraw.radius_exp}, \path{Szkraw.radius_exp_asymp}, \path{Szkraw.radius_no_uniform}\\[3pt]
The two examples in the paragraph ``How sharp is log-concavity?'' ($\mu=\frac1{10}$ and $\mu=\frac1{100}$) & \path{Szkraw.lc_violation_holds}, \path{Szkraw.lc_violation_fails}\\[3pt]
The combinations $(1-\lambda)K_{N,r}+\lambda G_{N,q,r}$ in the paragraph ``How sharp is log-concavity?'' and the seven listed failures & \path{Szkraw.combo_shape}, \path{Szkraw.combo_fails_4_5_10}, \path{Szkraw.combo_fails_3_5_20}, \path{Szkraw.combo_fails_1_2_20}, \path{Szkraw.combo_fails_2_5_60}, \path{Szkraw.combo_fails_1_3_100}, \path{Szkraw.combo_fails_3_10_200}, \path{Szkraw.combo_fails_2_7_200}
\end{tabular}
\end{center}
The entry for Theorem~\ref{thm:lc} formalizes Lemmas~\ref{lem:lc-ineq} and~\ref{lem:detours}, the operator induction, and Lemma~\ref{lem:kernel-rec} in converse form: the solution of~\eqref{eq:kernel-rec} yields the linearization coefficients. The Lean proofs of Corollaries~\ref{cor:kraw} and~\ref{cor:parity} obtain positivity of the components from the geometric-increment criterion of~\cite{P3} rather than from Theorem~\ref{thm:lc}; the statements are the ones given here. The Lean proof of Lemma~\ref{lem:corner-value} uses the algebra of~\cite{P3}.

Not formalized: the float and high-precision searches for the threshold $V_N$ in the paragraph ``How sharp is log-concavity?'', which are numerical evidence rather than exact claims, and the open problems of Section~\ref{sec:open}.

\begin{thebibliography}{9}

\bibitem{askey} R. Askey, \emph{Three notes on orthogonal polynomials}, Stichting Mathematisch Centrum, Afdeling Toegepaste Wiskunde, report TW~114/69, Amsterdam, October 1969; Part~II: \emph{Linearization of the product of orthogonal polynomials II}, \url{https://ir.cwi.nl/pub/7741}.
\bibitem{eagleson} G. K. Eagleson, \emph{A characterization theorem for positive definite sequences on the Krawtchouk polynomials}, Australian Journal of Statistics \textbf{11} (1969), no.~1, 29--38, \href{https://doi.org/10.1111/j.1467-842X.1969.tb00004.x}{\nolinkurl{doi:10.1111/j.1467-842X.1969.tb00004.x}}.
\bibitem{koelink} H. T. Koelink, \emph{$q$-Krawtchouk polynomials as spherical functions on the Hecke algebra of type B}, Transactions of the American Mathematical Society \textbf{352} (2000), no.~10, 4789--4813, \href{https://doi.org/10.1090/S0002-9947-00-02588-5}{\nolinkurl{doi:10.1090/S0002-9947-00-02588-5}}; preprint arXiv:q-alg/9605040v1 (1996), \url{https://arxiv.org/abs/q-alg/9605040v1}. Theorem and equation locators refer to the preprint.
\bibitem{miclo} L. Miclo, \emph{On the hypergroup property}, Annales de la Facult\'e des Sciences de Toulouse, s\'erie~6, \textbf{26} (2017), no.~2, 417--435, \href{https://doi.org/10.5802/afst.1539}{\nolinkurl{doi:10.5802/afst.1539}}.
\bibitem{szwarc} R. Szwarc, \emph{Orthogonal polynomials and Banach algebras}, in \emph{Inzell Lectures on Orthogonal Polynomials}, W. zu Castell, F. Filbir, and B. Forster (eds.), Advances in the Theory of Special Functions and Orthogonal Polynomials, vol.~2, Nova Science Publishers, 2005, 103--139; \S5, Problem~2. Author version (2004, 97--133; Problem~2 on printed p.~125): \url{https://www.math.uni.wroc.pl/~szwarc/pdf/inzell.pdf}.
\bibitem{szwarc2003} R. Szwarc, \emph{A necessary and sufficient condition for nonnegative product linearization of orthogonal polynomials}, Constructive Approximation \textbf{19} (2003), 565--573, \href{https://doi.org/10.1007/s00365-002-0524-z}{\nolinkurl{doi:10.1007/s00365-002-0524-z}}.
\bibitem{P2} L. P. Polzer, \emph{Increasing ratios in finite nonnegative linearization}, companion paper (2026).
\bibitem{P3} L. P. Polzer, \emph{Exact and perturbative criteria for finite nonnegative linearization}, companion paper (2026).
\bibitem{code} L. P. Polzer, \emph{Companion code for finite nonnegative linearization} (2026), revision \texttt{a25bf9a} (full identifier in the URL), \url{https://github.com/skypher/szkraw/tree/a25bf9a148047a6d909ae752e1ed1f8452aea394}.
\bibitem{docs} L. P. Polzer, \emph{Companion code for finite nonnegative linearization}, revision \texttt{b4c40c5} with the split documentation \path{papers/FORMALIZATION.md}, \path{papers/DIAGNOSTICS.md} and the fresh-clone record \path{runs/verify/RESULTS.md} (full identifier in the URL), \url{https://github.com/skypher/szkraw/tree/b4c40c54ec479a1a2e3df2739df35124f36f4673}.
\end{thebibliography}
\end{document}
